\section{Method}
\label{sec:method}

\subsection{Setting}

A standard run produces a checkpoint set
$\mathcal{C}=\{(\mathbf{w}_t,\, c_t)\}_{t=1}^{T}$, where $\mathbf{w}_t$ are the
parameters after epoch $t$ and $c_t$ is the scalar validation criterion the
pipeline already records at that epoch, on a split disjoint from the final
evaluation set. Pipelines differ in whether that criterion is minimized (a
loss) or maximized (a performance score), so we first remove the distinction.
Let $\delta=-1$ if the pipeline minimizes $c_t$ and $\delta=+1$ if it maximizes
it, and define the \emph{selection score}
\begin{equation}
  s_t \;=\; \delta \cdot c_t ,
  \label{eq:score}
\end{equation}
which is oriented so that larger is always better. Everything below is written
in terms of $s_t$ and is therefore identical for both kinds of pipeline.
Conventional selection keeps the single best state,
\begin{equation}
  t^{*}=\operatorname*{arg\,max}_{1\le t\le T} s_t ,
  \qquad
  \mathbf{w}_{\text{base}}=\mathbf{w}_{t^{*}} .
  \label{eq:baseline}
\end{equation}

\subsection{Criterion-Guided Top-$k$ Averaging}

Instead of the single maximizer we take the $k$ best-scoring epochs,
\begin{equation}
  \mathcal{S}_k=\operatorname*{arg\,max}_{\substack{\mathcal{S}\subseteq\{1,\dots,T\}\\ |\mathcal{S}|=k}}\ \sum_{t\in\mathcal{S}} s_t ,
  \label{eq:topk}
\end{equation}
that is, the indices of the $k$ largest $s_t$, and average their trainable
parameters element-wise:
\begin{equation}
  \mathbf{w}_{\text{avg}}=\frac{1}{|\mathcal{S}_k|}\sum_{t\in\mathcal{S}_k}\mathbf{w}_t .
  \label{eq:avg}
\end{equation}
Ties are broken in favour of later epochs, matching the convention by which
standard pipelines overwrite their best checkpoint on an equal score. Because
\cref{eq:topk} depends on $s_t$ only through its ordering, any strictly
monotone reparameterization of the criterion leaves $\mathcal{S}_k$ unchanged;
$\delta$ in \cref{eq:score} is the only task-dependent quantity in the method.

Two structural properties follow. First, $\mathcal{S}_1=\{t^{*}\}$, so
\cref{eq:baseline} is the $k{=}1$ case of \cref{eq:topk}: the baseline is
nested inside the method rather than being a separate procedure. Second,
$\mathcal{S}_{k-1}\subset\mathcal{S}_k$, so increasing $k$ only admits one
further checkpoint and never displaces an existing one.

\paragraph{Normalization statistics.}
Averaging is applied to trainable parameters only. Batch-Normalization
\cite{ioffe2015batch} running statistics are population estimates tied to the
particular weights that produced them, so those inherited from any single
member need not match $\mathbf{w}_{\text{avg}}$. Following SWA
\cite{izmailov2019averagingweightsleadswider} we reset them and recompute them
with forward passes over training data in training mode, without gradients or
optimizer updates. This is the method's only compute overhead and it is paid
once.

\subsection{Instantiating the Criterion}

The rule touches the task only through the pair $(c_t,\delta)$, which is why it
transfers without modification. We use whichever criterion the pipeline already
computes to select its own best checkpoint, so the method never introduces a
quantity the pipeline was not already tracking:

For \textbf{classification}, $c_t$ is the validation cross-entropy loss and
$\delta={-}1$, giving $s_t=-\ell^{\text{val}}_t$. For \textbf{detection and
segmentation}, $c_t$ is the validation fitness $\varphi_t$ and $\delta={+}1$,
giving $s_t=\varphi_t$; we read $\varphi_t$ straight from the Ultralytics
trainer \cite{jocher2023ultralytics}, which defines it as
$0.1\,\mathrm{mAP}_{50}+0.9\,\mathrm{mAP}_{50:95}$ for detection and as the sum
of that quantity over the box and mask components for segmentation. In both
cases it is the same scalar the pipeline uses for its own best checkpoint and
for early stopping, so baseline and rule are ranked by an identical signal.

No part of \cref{eq:topk,eq:avg} is task-specific, and we tune nothing per task
beyond setting $\delta$. The rule consequently inherits whatever biases the
pipeline's criterion already has: if the criterion is a poor proxy for test
performance, ranking by it will not repair that.

\subsection{Procedure and Cost}

\begin{algorithm}[t]
\caption{Criterion-guided top-$k$ checkpoint averaging}
\label{alg:topk}
\begin{algorithmic}[1]
\Require completed checkpoint set $\{(\mathbf{w}_t,c_t)\}_{t=1}^{T}$;
aggregation size $k\le T$; criterion direction
$\delta\in\{-1,+1\}$ ($-1$ if $c_t$ is a loss, $+1$ if it is a
score); training data $\mathcal{D}_{\text{train}}$
\Ensure final model $\mathcal{M}_{\text{avg}}$
\Statex \textbf{Rank} --- orient the criterion, then select
\For{$t=1,\dots,T$}
  \State $s_t \gets \delta\cdot c_t$ \Comment{larger is better}
\EndFor
\State $\mathcal{S}_k \gets$ indices of the $k$ largest $s_t$,
ordering ties by larger $t$
\Statex \textbf{Average} --- trainable parameters only
\State $\mathbf{w}_{\text{avg}} \gets \frac{1}{k}\sum_{t\in\mathcal{S}_k}\mathbf{w}_t$
\State load $\mathbf{w}_{\text{avg}}$ into $\mathcal{M}_{\text{avg}}$
\Statex \textbf{Recalibrate} --- only if normalization buffers exist
\If{$\mathcal{M}_{\text{avg}}$ has BatchNorm layers}
  \State reset their running mean and variance
  \State \textbf{for} a fixed number of batches from
  $\mathcal{D}_{\text{train}}$ \textbf{do} forward pass in training
  mode, no gradients, no optimizer step
\EndIf
\State \Return $\mathcal{M}_{\text{avg}}$
\end{algorithmic}
\end{algorithm}

\Cref{alg:topk} summarizes the procedure; setting $k{=}1$ recovers the baseline
line for line, which is what makes the two arms comparable. Training cost is
unchanged by construction, post-training cost is $k$ tensor additions plus one
recalibration pass, and inference cost is that of a single model. Storage is
the only real cost and it is bounded: because \cref{eq:topk} depends only on
the ordering of $s_t$, a running set of the $k$ best checkpoints seen so far
yields the same $\mathcal{S}_k$ as ranking the full set offline, so $k{+}1$
states suffice at any moment. Our detection and segmentation runs implement
exactly this, pruning a new checkpoint immediately if it falls outside the
current top $k$; the classification runs retain all epochs and we report the
counts in \cref{sec:selection}.

\paragraph{Implementation details.}
Two choices keep the comparison clean. The Ultralytics trainer validates
through an exponential moving average of the weights by default; we disable
that smoothing so the recorded fitness, the early-stopping signal and the
stored weights all refer to the same raw parameters, since otherwise
checkpoints would be ranked by a quantity computed from weights other than
those being averaged. We also take the baseline from the top-ranked entry of
the same table our rule ranks rather than from the framework's exported best
checkpoint, which is serialized at half precision: comparing it against a
full-precision average would confound the selection rule with a precision
change.
